% Kickoff deck: four exploratory findings on the day-ahead forecast for NHS sites.
% Build with build.sh (pdfLaTeX, two passes). Text follows ASD-STE100 Simplified Technical English.
% Look: the Oxford / SENSE slide template (Oxford blue, serif titles, Helvetica text).
% Slide size 160 x 90 mm. Figures come from code/scripts/deck/ at their exact size on the slide.
\documentclass[aspectratio=169,9pt]{beamer}
\usepackage[T1]{fontenc}
\usepackage{tgheros}                 % Helvetica (TeX Gyre Heros) for text
\usepackage[tracking=true]{microtype}
\usepackage{tikz}
\usepackage{graphicx}
\usefonttheme{professionalfonts}
\setbeamertemplate{navigation symbols}{}
\setbeamersize{text margin left=0mm,text margin right=0mm}

% ---- event line on the title slide: set the date when it is fixed ---------------------------
\newcommand{\deckevent}{SENSE healthcare kickoff}
\def\deckdate{}
% ---- backup slides: the three earlier findings (heat pumps, hot days, verification) --------
% \backuptrue adds them after the discussion slide
\newif\ifbackup
\backupfalse

% ---- fonts: Charter for titles and numbers, Helvetica for text --------------------------------
\newcommand{\titlefont}{\fontfamily{bch}\selectfont}

% ---- colours of the template ----------------------------------------------------------------
\definecolor{oxblue}{HTML}{002147}
\definecolor{shell}{HTML}{F1EEE9}
\definecolor{charcoal}{HTML}{211D1C}
\definecolor{ash}{HTML}{61615F}
\definecolor{stone}{HTML}{D9D8D6}
\setbeamercolor{normal text}{fg=charcoal,bg=white}

% ---- placement: millimetres from the top-left corner of the slide ----------------------------
\newcommand{\at}[3]{\node[anchor=north west,inner sep=0pt,outer sep=0pt] at ([xshift=#1mm,yshift=-#2mm]current page.north west) {#3};}
\newcommand{\atr}[3]{\node[anchor=north east,inner sep=0pt,outer sep=0pt] at ([xshift=#1mm,yshift=-#2mm]current page.north west) {#3};}
\newcommand{\fig}[3]{\at{#1}{#2}{\includegraphics{figures/#3.pdf}}}

\newcommand{\lab}[1]{{\bfseries\fontsize{5.4}{7}\selectfont\color{oxblue}\textls[90]{\MakeUppercase{#1}}}}
\newcommand{\bodytxt}{\fontsize{7.3}{9.6}\selectfont\color{charcoal}}
\newcommand{\slidehead}[3]{%
  \at{7.2}{6.4}{\lab{#1}}
  \at{7.2}{9.7}{\parbox[t]{146mm}{\raggedright\titlefont\fontsize{14.6}{17}\selectfont\color{oxblue}#2}}
  \at{7.2}{17.9}{\parbox[t]{146mm}{\raggedright\fontsize{7.2}{9.4}\selectfont\color{ash}#3}}}
\newcommand{\slidefoot}{%
  \at{7.2}{84.3}{{\fontsize{4.6}{6}\selectfont\color{ash}Saïd Business School, University of Oxford \enspace·\enspace Smart Energy Data Service (SENSE)}}
  \atr{152.8}{84.0}{{\bfseries\fontsize{5.4}{6}\selectfont\color{oxblue}\insertframenumber}}}
\newcommand{\bignum}[5]{% x, y, width, number, text
  \at{#1}{#2}{{\titlefont\fontsize{16}{17}\selectfont\color{oxblue}#4}}
  \node[anchor=north west,inner sep=0pt,outer sep=0pt] at ([xshift=#1mm,yshift=-#2mm-7.4mm]current page.north west) {\parbox[t]{#3mm}{\raggedright\fontsize{6.2}{7.9}\selectfont\color{ash}#5}};}
% the three text blocks of a finding slide
\newcommand{\blk}[2]{\lab{#1}\par\vspace{1mm}{\bodytxt\raggedright #2\par}\vspace{2.5mm}}
\newcommand{\hyp}[2]{\par\hangindent=5mm\hangafter=1\noindent\makebox[5mm][l]{\bfseries\color{oxblue}#1}#2\par\vspace{0.9mm}}
\newcommand{\textcol}[1]{%
  \draw[stone,line width=0.35pt] ([xshift=100.4mm,yshift=-24.5mm]current page.north west) -- ++(0,-55.5mm);
  \at{103.6}{24.0}{\parbox[t]{49.2mm}{#1}}}
\newcommand{\bul}[1]{\par\hangindent=3.1mm\hangafter=1\noindent\makebox[3.1mm][l]{\raisebox{0.45mm}{\color{oxblue}\rule{0.95mm}{0.95mm}}}#1\par\vspace{1.5mm}}
\newcommand{\boxcol}[3]{% x, label, content
  \fill[shell] ([xshift=#1mm,yshift=-24mm]current page.north west) rectangle ++(45.4mm,-55.5mm);
  \node[anchor=north west,inner sep=0pt,outer sep=0pt] at ([xshift=#1mm+4mm,yshift=-28mm]current page.north west) {\parbox[t]{37.6mm}{{\bfseries\fontsize{8.6}{10}\selectfont\color{oxblue}#2}\par\vspace{2.6mm}{\fontsize{7.2}{9.6}\selectfont\color{charcoal}\raggedright #3}}};}
% a logo appears when its file is in logos/ (pdf or png); without the file the slot stays empty
\newcommand{\putlogo}[3]{\IfFileExists{logos/#1.pdf}{\at{#2}{#3}{\includegraphics[height=9mm]{logos/#1.pdf}}}{\IfFileExists{logos/#1.png}{\at{#2}{#3}{\includegraphics[height=9mm]{logos/#1.png}}}{}}}

\begin{document}

% =============================================================================================
% 1  Title
\begin{frame}[plain]\begin{tikzpicture}[remember picture,overlay]
  \fill[oxblue] ([xshift=98.4mm]current page.north west) rectangle (current page.south east);
  \fig{101.2}{5.5}{cover_sites}
  \at{104}{79.2}{\parbox[t]{50mm}{\raggedright\fontsize{5.2}{6.8}\selectfont\color{white}134 NHS sites with half-hourly electricity data.\\ Gold: sites with on-site generation.}}
  \at{7.2}{8.6}{\lab{Smart Energy Data Service (SENSE) \enspace·\enspace SDR UK}}
  \at{7.2}{27}{\parbox[t]{88mm}{\raggedright\titlefont\fontsize{19.5}{24}\selectfont\color{oxblue}NHS site electricity demand:\\ four exploratory findings}}
  \at{7.2}{47.5}{\parbox[t]{84mm}{\raggedright\fontsize{9}{12.2}\selectfont\color{ash}What a day-ahead forecast can do, and what limits it. Half-hourly data for 134 NHS sites, 2022 to 2026}}
  \at{7.2}{60.5}{{\bfseries\fontsize{8.4}{10}\selectfont\color{oxblue}Professor Siddharth Arora \enspace·\enspace Nicolas Malz}}
  \at{7.2}{65.5}{{\fontsize{6.4}{8}\selectfont\color{ash}Saïd Business School, University of Oxford \enspace·\enspace \deckevent\ifx\deckdate\empty\else\enspace·\enspace\deckdate\fi}}
  \putlogo{oxford_logo}{7.2}{74}
  \putlogo{sbs_logo}{24}{74}
  \putlogo{sense_logo}{48}{74}
\end{tikzpicture}\end{frame}

% =============================================================================================
% 2  Finding 1: the forecast and its value
\begin{frame}[plain]\begin{tikzpicture}[remember picture,overlay]
  \slidehead{Exploratory finding 1 \enspace·\enspace forecast value}
       {Zero-shot models lead, but get only half of the possible value}
       {Day-ahead forecasts for 91 sites on 69 test days. Value: the peak cut of a battery that plans on the forecast.}
  \fig{7.2}{24.6}{forecast}
  \bignum{29}{66}{28}{11.5 \%}{error of the best model, with no site training}
  \bignum{65.5}{66}{30}{53 \%}{of the possible peak cut, with the best forecast}
  \textcol{%
    \blk{Implication}{A model with no site training is the baseline to beat. Two weeks of history are sufficient. A perfect daily level increases the value to 63 \%. A new plan at 06:00 adds 2 points.}
    \blk{Research avenue}{Train the model for the decision, not for the mean error. Find the time and the height of the peak, not the level of the day.}
    \blk{Hypotheses}{\hyp{H1}{A model that is trained for the peak cut gets more value with the same error.}\hyp{H2}{The best quantile for the plan is different for each site.}}}
  \slidefoot
\end{tikzpicture}\end{frame}

% =============================================================================================
% 3  Finding 2: on-site generation
\begin{frame}[plain]\begin{tikzpicture}[remember picture,overlay]
  \slidehead{Exploratory finding 2 \enspace·\enspace on-site generation}
       {Sites with on-site generation have 79 \% of the forecast error}
       {12 of the 91 sites. Map: forecast error of each site in kWh. Chart: grid import of one site on a day when its plant stops.}
  \fig{7.2}{24.2}{error_map}
  \fig{55.6}{24.6}{switch_day}
  \bignum{58.4}{66}{38}{43 \%}{of all error is on 2 \% of the site-days: the days with a plant start or stop}
  \textcol{%
    \blk{Implication}{The meter shows demand minus generation. A plant start or stop changes the import by more than its mean in two hours. Size explains half of the error at these sites, the plant the other half.}
    \blk{Research avenue}{Get plant schedules or generation meter data from two or three trusts. Then forecast demand and generation as two series.}
    \blk{Hypotheses}{\hyp{H1}{With the plant schedule, the error at these sites decreases by half.}\hyp{H2}{The site knows most plant starts and stops one day before.}}}
  \slidefoot
\end{tikzpicture}\end{frame}

% =============================================================================================
% 4  Finding 3: patient activity
\begin{frame}[plain]\begin{tikzpicture}[remember picture,overlay]
  \slidehead{Exploratory finding 3 \enspace·\enspace patient activity}
       {Strikes and visitor numbers do not change electricity demand}
       {Left: demand on 34 doctor strike days against normal days, 132 sites. Right: daily visitors against daily demand at 2 sites, 2024.}
  \fig{7.2}{24.6}{strike_profile}
  \fig{55.6}{24.6}{footfall}
  \bignum{12}{66}{38}{$-$0.4 \%}{demand from 08:00 to 18:00 on strike days (95 \% interval: $-$1.6 to +0.8 \%)}
  \bignum{60.4}{66}{37}{0.00}{change in demand for 1 \% more visitors. Mobility data: The Oxford Partnership}
  \textcol{%
    \blk{Implication}{On strike days, hospitals cancelled operations and appointments. Demand did not change. A model that gets the visitor count of the day is not better. Buildings and plant set the demand, not patients.}
    \blk{Research avenue}{Do the same test with daily bed occupancy from the trusts. Do the test for gas, because heat can follow occupancy.}
    \blk{Hypotheses}{\hyp{H1}{Daily bed occupancy explains less than 1 \% of the daily change in electricity.}\hyp{H2}{Activity changes demand only when a department closes fully.}}}
  \slidefoot
\end{tikzpicture}\end{frame}

% =============================================================================================
% 5  Finding 4: heat waves and cold waves
\begin{frame}[plain]\begin{tikzpicture}[remember picture,overlay]
  \slidehead{Exploratory finding 4 \enspace·\enspace heat and cold waves}
       {On heat-wave days the forecast loses a third of its value}
       {91 sites, 342 test days, forecast with no weather input. Wave day: site temperature beyond the 5th or 95th percentile.}
  \fig{7.2}{24.6}{waves}
  \bignum{29}{66}{28}{+1.5 points}{of error on heat-wave days: 14.7 \% against 13.2 \%}
  \bignum{65.5}{66}{30}{33 \%}{of the possible peak cut on heat-wave days. Other days: 50 \%}
  \textcol{%
    \blk{Implication}{The error shows little, the decision shows much. Demand is above the forecast, the plan is too low, and the battery is empty before the peak. With a weather forecast as input, the best model gets 40 \%.}
    \blk{Research avenue}{Score forecasts for rare days in decisions. Test synthetic heat-wave days (Lu et al., 2026) and a plan on a higher quantile.}
    \blk{Hypotheses}{\hyp{H1}{Sites with cooling plant lose the most value on heat-wave days.}\hyp{H2}{Synthetic heat-wave days recover value at sites that follow the temperature.}}}
  \slidefoot
\end{tikzpicture}\end{frame}

% =============================================================================================
% 6  Discussion
\begin{frame}[plain]\begin{tikzpicture}[remember picture,overlay]
  \slidehead{Discussion}
       {Facts, open points and questions for this session}
       {All results are exploratory: one data panel, one test period. No trust has checked the results.}
  \boxcol{7.2}{Facts}{%
    \bul{134 sites in 24 trusts. Half-hourly electricity, December 2022 to May 2026.}
    \bul{Forecast test: 91 sites, 10 models, 69 to 342 test days.}
    \bul{12 of the 91 sites have on-site generation. Their meters show grid import only.}
    \bul{Public load data: the zero-shot model is level with a method from 2026.}}
  \boxcol{57.3}{Not resolved}{%
    \bul{Plant schedules and generation meter data: we have none.}
    \bul{Daily bed occupancy: public data are monthly or quarterly.}
    \bul{Value is one decision here: the peak cut of a battery. Other decisions can give other results.}
    \bul{Can we name sites in publications?}}
  \boxcol{107.4}{Questions for this session}{%
    \bul{Which trusts can supply plant schedules or generation meter data?}
    \bul{Which group has daily occupancy or activity data that we can link?}
    \bul{Which decisions do estates teams make with a forecast, and at what time?}
    \bul{Can all groups use one site identifier and one test protocol?}}
  \slidefoot
\end{tikzpicture}\end{frame}

% =============================================================================================
% Backup slides: the three earlier findings. Set \backuptrue at the top to add them.
\ifbackup
% Backup: heat electrification against the grid connection
\begin{frame}[plain]\begin{tikzpicture}[remember picture,overlay]
  \slidehead{Earlier result \enspace·\enspace heat electrification}
       {Heat pumps put half of the sites above their grid connection}
       {Scenario: heat pumps replace all gas heat. Map: 631 sites in England. Chart: 31 panel sites in the weather of 66 years.}
  \fig{7.2}{24.2}{headroom_map}
  \fig{55.6}{24.6}{headroom_winters}
  \bignum{58.4}{66}{38}{48 \%}{of 631 sites need more power than the connection can supply}
  \textcol{%
    \blk{Implication}{The winter peak is the limit, not the annual energy. Mild winters and cold winters give almost the same result. Storage keeps 5 of 14 sites below the limit.}
    \blk{Research avenue}{Add network capacity data. Find the lowest-cost mix of storage, insulation and a larger connection for each site.}
    \blk{Hypotheses}{\hyp{H1}{Gas for the coldest hours keeps most sites below the limit.}\hyp{H2}{Sites near the limit need less than 2 hours of storage.}}}
  \slidefoot
\end{tikzpicture}\end{frame}

% Backup: the response to hot days
\begin{frame}[plain]\begin{tikzpicture}[remember picture,overlay]
  \slidehead{Earlier result \enspace·\enspace weather response}
       {Hot days increase demand at some sites and decrease it at others}
       {83 sites without on-site generation. On hot days, 33 sites use more electricity and 21 sites use less.}
  \fig{7.2}{24.2}{heat_map}
  \fig{55.6}{24.6}{heat_explained}
  \bignum{58.4}{66}{38}{45 \%}{of the difference between sites comes from one meter pattern. Registers explain 0 \%.}
  \textcol{%
    \blk{Implication}{One national heat coefficient is not correct for a site. Building registers do not explain the difference between sites. The meter does.}
    \blk{Research avenue}{Get plant lists from two or three trusts. Then find cooling plant and heating plant at all sites from the meter pattern.}
    \blk{Hypotheses}{\hyp{H1}{Sites that use more on hot days have cooling plant.}\hyp{H2}{Sites that use less have heating that operates in summer.}}}
  \slidefoot
\end{tikzpicture}\end{frame}

% Backup: verification from the meter
\begin{frame}[plain]\begin{tikzpicture}[remember picture,overlay]
  \slidehead{Earlier result \enspace·\enspace verification}
       {The meter shows when sites install solar panels or heat pumps}
       {Left: how demand follows sunshine at 107 sites. Right: funded solar sites report 84 \% more generation; other sites 39 \%.}
  \fig{7.2}{24.2}{solar_lines}
  \fig{51.4}{24.2}{solar_programme_map}
  \textcol{%
    \blk{Implication}{Half-hourly meters can check public programmes at each site. The meter gives the date and the size of an installation.}
    \blk{Research avenue}{Use funded solar sites and new heat pumps as natural experiments. Calculate the cost of each kWh that a programme saves.}
    \blk{Hypotheses}{\hyp{H1}{After one full year, heat pump sites use 8 \% more electricity and 9 \% less gas.}\hyp{H2}{Sites with a large array export most of the solar energy.}}}
  \slidefoot
\end{tikzpicture}\end{frame}

\fi

\end{document}
